\begin{abstract}
    % Diffusion models have emerged as highly effective generative models; however, their considerable computational demands present a significant obstacle to on-device deployment. In prior work, Modulated Diffusion (MoDiff)~\cite{gao2025modulateddiffusionacceleratinggenerative} introduced a framework that offers an innovative and practical acceleration strategy based on modulated quantization and error compensation. The benefits and optimizations of this approach are supported by rigorous theoretical analysis. Simulated experiments further indicate that reducing activation quantization precision from 8 bits to 3 bits does not result in observable performance degradation.

    % In this study, we present an engineering-oriented extension of MoDiff aimed at constructing a more robust and efficient simulation and evaluation framework. The original MoDiff implementation exhibited limited capacity to demonstrate substantial end-to-end speedups, motivating us to investigate architectural and implementation-level enhancements. To address this limitation, we incorporated deployable and benchmarkable components into the framework, thereby enabling systematic and reproducible assessment of performance improvements. Our implementation employs fused int8 and int4 kernels built on CUTLASS (CUDA Templates for Linear Algebra Subroutines and Solvers), which yield substantial runtime performance gains. In addition, we conducted a comprehensive and systematic analysis of the model’s runtime characteristics. This detailed investigation identified key performance bottlenecks and directly informed the subsequent optimization strategies.
    Diffusion models have demonstrated remarkable effectiveness as generative models, but their high computational cost, caused by the iterative denoising process and computationally heavy backbone networks, remains a major barrier to on-device and real-time deployment. Modulated Diffusion (MoDiff) provides a post-training quantization framework that supports highly efficient low-bit activation quantization. However, prior work evaluates the acceleration benefits of MoDiff only through operation-count analysis, without validating them on real hardware. In this paper, we present a kernel implementation of MoDiff for practical hardware acceleration. In particular, we propose a cache-update fusion paradigm that eliminates additional I/O operations during inference. Experiments and ablation studies demonstrate that our implementation achieves up to a 1.8× runtime speedup over FP32 models and up to a 42.2\% reduction in memory I/O usage compared to FP32.
    %\red{add memory saving results}.
\end{abstract}